\documentclass[12pt, a4paper]{article}

\usepackage[utf8]{inputenc}
\usepackage[T1]{fontenc}
\usepackage[brazil]{babel}
\usepackage{amsmath, amssymb}
\usepackage{geometry}
\usepackage{listings}
\usepackage{xcolor}
\usepackage{enumitem}
\usepackage{mdframed}
\usepackage{tikz}
\usetikzlibrary{trees, arrows.meta}

\geometry{margin=2.5cm}

% --- Estilo dos blocos de código ---
\lstdefinestyle{cstyle}{
    language=C,
    basicstyle=\ttfamily\small,
    keywordstyle=\color{blue}\bfseries,
    commentstyle=\color{gray}\itshape,
    stringstyle=\color{orange},
    numbers=left,
    numberstyle=\tiny\color{gray},
    stepnumber=1,
    frame=single,
    backgroundcolor=\color{gray!8},
    breaklines=true,
    tabsize=4,
}

\lstdefinestyle{javastyle}{
    language=Java,
    basicstyle=\ttfamily\small,
    keywordstyle=\color{purple}\bfseries,
    commentstyle=\color{gray}\itshape,
    stringstyle=\color{orange},
    numbers=left,
    numberstyle=\tiny\color{gray},
    stepnumber=1,
    frame=single,
    backgroundcolor=\color{purple!5},
    breaklines=true,
    tabsize=4,
}

% --- Caixa de resposta ---
\newenvironment{resposta}{%
    \begin{mdframed}[linecolor=black!30, backgroundcolor=white,
                     innertopmargin=10pt, innerbottommargin=30pt,
                     skipabove=6pt]
    \textbf{Resposta:}\\[4pt]
}{%
    \end{mdframed}
}

% --- Estilo dos nós da árvore ---
\tikzset{
    no/.style={circle, draw=black, thick, minimum size=8mm, inner sep=1pt},
    nulo/.style={draw=none, fill=none},
}

% -------------------------------------------------------
\title{\textbf{Lista de Exercícios}\\[6pt]
       \large Árvore Binária e Árvore Binária de Busca}
\author{Lucas Rafael Pessoa do Nascimento}
\date{}

\begin{document}

\maketitle

\begin{center}
    \textit{Os exercícios cobrem implementação, percurso, análise de complexidade\\
    e propriedades de árvores binárias e árvores binárias de busca (ABB).\\
    Utilize as structs/classes apresentadas em cada seção como base.}
\end{center}

\bigskip

% =========================================================
\section*{Estruturas de Referência}
% =========================================================

\textbf{Em C:}
\begin{lstlisting}[style=cstyle]
typedef struct No {
    int valor;
    struct No *esq, *dir;
} No;

No* novoNo(int v) {
    No *n = (No*) malloc(sizeof(No));
    n->valor = v;
    n->esq = n->dir = NULL;
    return n;
}
\end{lstlisting}

\textbf{Em Java:}
\begin{lstlisting}[style=javastyle]
class No {
    int valor;
    No esq, dir;
    No(int v) { valor = v; esq = dir = null; }
}
\end{lstlisting}

\bigskip

% =========================================================
\section*{Parte I — Conceitos Básicos e Percursos}
% =========================================================

% ---------------------------------------------------------
\subsection*{Exercício 1 — Contar nós (recursivo em C)}

\begin{lstlisting}[style=cstyle]
int contarNos(No *raiz) {
    if (raiz == NULL) return 0;
    return 1 + contarNos(raiz->esq) + contarNos(raiz->dir);
}
\end{lstlisting}

\begin{enumerate}[label=(\alph*)]
    \item Explique o caso base e o passo recursivo.
    \item Qual é a complexidade de tempo em função do número de nós $n$? Justifique com um somatório.
    \item Reescreva a função de forma \textbf{iterativa} usando uma pilha (ou fila).
\end{enumerate}

\begin{resposta}
\vspace{4cm}
\end{resposta}

% ---------------------------------------------------------
\subsection*{Exercício 2 — Altura da árvore (recursivo em Java)}

\begin{lstlisting}[style=javastyle]
int altura(No raiz) {
    if (raiz == null) return -1;
    return 1 + Math.max(altura(raiz.esq), altura(raiz.dir));
}
\end{lstlisting}

\begin{enumerate}[label=(\alph*)]
    \item Por que o caso base retorna $-1$ (e não $0$)?
    \item Para a árvore abaixo, calcule manualmente a altura:

\begin{center}
\begin{tikzpicture}[level distance=1.2cm,
    level 1/.style={sibling distance=4cm},
    level 2/.style={sibling distance=2cm},
    level 3/.style={sibling distance=1.2cm}]
  \node[no]{10}
    child { node[no]{5}
        child { node[no]{3}
            child { node[no]{1} }
            child[missing]{}
        }
        child { node[no]{7} }
    }
    child { node[no]{15}
        child[missing]{}
        child { node[no]{20} }
    };
\end{tikzpicture}
\end{center}

    \item No \textbf{pior caso} (árvore degenerada/lista), qual é a altura? E no \textbf{melhor caso} (árvore completa com $n$ nós)?
    \item Qual a complexidade de tempo de \texttt{altura()}?
\end{enumerate}

\begin{resposta}
\vspace{4cm}
\end{resposta}

% ---------------------------------------------------------
\subsection*{Exercício 3 — Percursos: Pré-ordem, In-ordem e Pós-ordem}

Dada a árvore do Exercício 2, escreva a sequência de valores visitados em cada percurso:

\begin{enumerate}[label=(\alph*)]
    \item \textbf{Pré-ordem} (raiz, esquerda, direita)
    \item \textbf{In-ordem} (esquerda, raiz, direita)
    \item \textbf{Pós-ordem} (esquerda, direita, raiz)
    \item Implemente o percurso \textbf{in-ordem} em Java de forma \textbf{iterativa} usando uma pilha (\texttt{Stack<No>}).
    \item Qual percurso produz os valores em ordem crescente em uma ABB? Por quê?
\end{enumerate}

\begin{resposta}
\vspace{5cm}
\end{resposta}

\newpage
% ---------------------------------------------------------
\subsection*{Exercício 4 — Contar folhas (C)}

\begin{lstlisting}[style=cstyle]
int contarFolhas(No *raiz) {
    if (raiz == NULL) return 0;
    if (raiz->esq == NULL && raiz->dir == NULL) return 1;
    return contarFolhas(raiz->esq) + contarFolhas(raiz->dir);
}
\end{lstlisting}

\begin{enumerate}[label=(\alph*)]
    \item Quantas folhas tem a árvore do Exercício 2? Verifique com a função.
    \item Uma árvore completa com $k$ níveis tem quantas folhas? Expresse como potência de 2.
    \item Escreva uma função \texttt{contarInternosC} que conta apenas os nós internos (não-folhas).
\end{enumerate}

\begin{resposta}
\vspace{4cm}
\end{resposta}

% =========================================================
\section*{Parte II — Árvore Binária de Busca (ABB)}
% =========================================================

% ---------------------------------------------------------
\subsection*{Exercício 5 — Busca em ABB (Java)}

\begin{lstlisting}[style=javastyle]
boolean buscar(No raiz, int x) {
    if (raiz == null) return false;
    if (x == raiz.valor) return true;
    if (x < raiz.valor) return buscar(raiz.esq, x);
    return buscar(raiz.dir, x);
}
\end{lstlisting}

\begin{enumerate}[label=(\alph*)]
    \item Trace a busca por $7$ na árvore do Exercício 2. Quais nós são visitados?
    \item Qual é a complexidade da busca no \textbf{melhor caso}? E no \textbf{pior caso}?
    \item Reescreva \texttt{buscar} de forma \textbf{iterativa} (sem recursão).
    \item Por que a busca em ABB é mais eficiente que a busca linear em um vetor não ordenado?
\end{enumerate}

\begin{resposta}
\vspace{4.5cm}
\end{resposta}

% ---------------------------------------------------------
\subsection*{Exercício 6 — Inserção em ABB (C)}

\begin{lstlisting}[style=cstyle]
No* inserir(No *raiz, int v) {
    if (raiz == NULL) return novoNo(v);
    if (v < raiz->valor)
        raiz->esq = inserir(raiz->esq, v);
    else if (v > raiz->valor)
        raiz->dir = inserir(raiz->dir, v);
    return raiz;
}
\end{lstlisting}

\begin{enumerate}[label=(\alph*)]
    \item Insira os valores $\{10, 5, 15, 3, 7, 20, 1\}$ nessa ordem. Desenhe a ABB resultante.
    \item Agora insira $\{1, 3, 5, 7, 10, 15, 20\}$ nessa ordem. O que acontece com a árvore?
    \item Compare as alturas das duas árvores. O que isso implica para a complexidade das operações?
    \item A função acima ignora duplicatas. Modifique-a para \textbf{contar} quantas vezes cada valor foi inserido (armazene um campo \texttt{freq} no nó).
\end{enumerate}

\begin{resposta}
\vspace{5cm}
\end{resposta}

\newpage
% ---------------------------------------------------------
\subsection*{Exercício 7 — Mínimo e Máximo em ABB (Java)}

\begin{lstlisting}[style=javastyle]
int minimo(No raiz) {
    if (raiz.esq == null) return raiz.valor;
    return minimo(raiz.esq);
}

int maximo(No raiz) {
    if (raiz.dir == null) return raiz.valor;
    return maximo(raiz.dir);
}
\end{lstlisting}

\begin{enumerate}[label=(\alph*)]
    \item Explique por que o mínimo está sempre à esquerda e o máximo à direita.
    \item Reescreva ambas as funções de forma \textbf{iterativa}.
    \item Em uma ABB com $n$ nós balanceada, quantas comparações são feitas para encontrar o mínimo? E em uma ABB degenerada?
    \item O mínimo de uma ABB corresponde a qual posição no percurso in-ordem?
\end{enumerate}

\begin{resposta}
\vspace{4cm}
\end{resposta}

% ---------------------------------------------------------
\subsection*{Exercício 8 — Verificar se é ABB (C)}

A função abaixo verifica se uma árvore binária é uma ABB válida:

\begin{lstlisting}[style=cstyle]
int ehABB(No *raiz, int min, int max) {
    if (raiz == NULL) return 1;
    if (raiz->valor <= min || raiz->valor >= max) return 0;
    return ehABB(raiz->esq, min, raiz->valor) &&
           ehABB(raiz->dir, raiz->valor, max);
}
/* Chamada inicial: ehABB(raiz, INT_MIN, INT_MAX) */
\end{lstlisting}

\begin{enumerate}[label=(\alph*)]
    \item Por que simplesmente comparar \texttt{raiz->valor} com \texttt{raiz->esq->valor} e \texttt{raiz->dir->valor} não é suficiente? Dê um contraexemplo.
    \item Trace a execução para a árvore abaixo e determine se é uma ABB válida:

\begin{center}
\begin{tikzpicture}[level distance=1.2cm,
    level 1/.style={sibling distance=3.5cm},
    level 2/.style={sibling distance=1.8cm}]
  \node[no]{8}
    child { node[no]{3}
        child { node[no]{1} }
        child { node[no]{6} }
    }
    child { node[no]{12}
        child { node[no]{5} }
        child { node[no]{14} }
    };
\end{tikzpicture}
\end{center}

    \item Qual é a complexidade de tempo de \texttt{ehABB}?
\end{enumerate}

\begin{resposta}
\vspace{4.5cm}
\end{resposta}

% =========================================================
\section*{Parte III — Análise de Complexidade e Desafios}
% =========================================================

% ---------------------------------------------------------
\subsection*{Exercício 9 — Complexidade das operações em ABB}

Preencha a tabela considerando uma ABB com $n$ nós:

\bigskip
\begin{center}
\renewcommand{\arraystretch}{2.2}
\begin{tabular}{|l|c|c|c|}
\hline
\textbf{Operação} & \textbf{Melhor caso} & \textbf{Caso médio} & \textbf{Pior caso} \\
\hline
Busca        & & & \\
\hline
Inserção     & & & \\
\hline
Mínimo/Máximo & & & \\
\hline
Percurso in-ordem & & & \\
\hline
Contar nós   & & & \\
\hline
Altura       & & & \\
\hline
\end{tabular}
\end{center}

\bigskip
\begin{enumerate}[label=(\alph*)]
    \item Quando o pior caso de busca e inserção é $O(n)$? Descreva a forma da árvore nessa situação.
    \item O que é uma árvore \textbf{balanceada} e por que ela garante $O(\log n)$?
\end{enumerate}

\begin{resposta}
\vspace{3cm}
\end{resposta}

% ---------------------------------------------------------
\subsection*{Exercício 10 — Percurso por nível (BFS em Java)}

\begin{lstlisting}[style=javastyle]
void porNivel(No raiz) {
    if (raiz == null) return;
    Queue<No> fila = new LinkedList<>();
    fila.add(raiz);
    while (!fila.isEmpty()) {
        No atual = fila.poll();
        System.out.print(atual.valor + " ");
        if (atual.esq != null) fila.add(atual.esq);
        if (atual.dir != null) fila.add(atual.dir);
    }
}
\end{lstlisting}

\begin{enumerate}[label=(\alph*)]
    \item Trace a execução para a árvore do Exercício 2. Qual é a ordem de saída?
    \item Por que este algoritmo usa uma \textbf{fila} e não uma \textbf{pilha}? O que mudaria se usasse pilha?
    \item Qual é a complexidade de tempo e espaço de \texttt{porNivel}?
    \item Modifique \texttt{porNivel} para imprimir cada nível em uma linha separada.
\end{enumerate}

\begin{resposta}
\vspace{5cm}
\end{resposta}

% ---------------------------------------------------------
\subsection*{Exercício 11 — Remoção em ABB (C)}

A remoção em ABB tem três casos. Implemente a função:

\begin{lstlisting}[style=cstyle]
No* remover(No *raiz, int v) {
    if (raiz == NULL) return NULL;
    if (v < raiz->valor)
        raiz->esq = remover(raiz->esq, v);
    else if (v > raiz->valor)
        raiz->dir = remover(raiz->dir, v);
    else {
        /* Caso 1: folha */
        if (raiz->esq == NULL && raiz->dir == NULL) {
            free(raiz); return NULL;
        }
        /* Caso 2: um filho */
        if (raiz->esq == NULL) {
            No *tmp = raiz->dir; free(raiz); return tmp;
        }
        if (raiz->dir == NULL) {
            No *tmp = raiz->esq; free(raiz); return tmp;
        }
        /* Caso 3: dois filhos — substitui pelo minimo da subarvore direita */
        No *suc = raiz->dir;
        while (suc->esq != NULL) suc = suc->esq;
        raiz->valor = suc->valor;
        raiz->dir = remover(raiz->dir, suc->valor);
    }
    return raiz;
}
\end{lstlisting}

\begin{enumerate}[label=(\alph*)]
    \item Descreva os três casos de remoção com suas próprias palavras.
    \item Na árvore do Exercício 2, remova o nó $5$. Desenhe o resultado.
    \item Na árvore do Exercício 2, remova o nó $10$ (raiz). Qual nó assume o lugar da raiz? Desenhe.
    \item Qual é a complexidade da remoção no pior caso?
\end{enumerate}

\begin{resposta}
\vspace{5cm}
\end{resposta}

\newpage
% ---------------------------------------------------------
\subsection*{Exercício 12 — Desafio: Espelhamento e Simetria}

\begin{enumerate}[label=(\alph*)]
    \item Escreva uma função recursiva em Java \texttt{espelhar(No raiz)} que inverte uma árvore binária (troca todos os filhos esquerdos com os direitos).
    \item Escreva uma função \texttt{ehSimetrica(No raiz)} que retorna \texttt{true} se a árvore é espelho de si mesma (palíndromo estrutural). Use uma função auxiliar \texttt{espelhos(No a, No b)}.
    \item Aplique \texttt{espelhar} à árvore do Exercício 2 e desenhe o resultado.
    \item Determine a complexidade de \texttt{espelhar} e de \texttt{ehSimetrica}.
\end{enumerate}

\begin{resposta}
\vspace{6cm}
\end{resposta}

\end{document}
